\documentclass[../cours_2de.tex]{subfiles}
\begin{document}
\setcounter{chapter}{11}
\chapter{Fonctions : définir, représenter, modéliser}
\label{sec:ch12}
\begin{objectif}
Comprendre une fonction au-delà d'une formule ; déterminer son domaine ; utiliser courbes, tableaux et programmes en connaissant leurs limites.
\end{objectif}
\section{Une sortie pour chaque entrée}
\begin{definition}
Une fonction $f$ définie sur un ensemble $D$ associe à chaque $x\in D$ un unique réel $f(x)$, son image.
Un antécédent de $y$ est un $x\in D$ tel que $f(x)=y$. Une valeur peut avoir aucun, un ou plusieurs antécédents.
\end{definition}
Pour $f(x)=x^2$ sur $\R$, $f(-3)=9$ ; $9$ possède les antécédents $-3$ et $3$, $0$ en possède un, et $-1$ aucun.
L'écriture $f(x)$ ne signifie pas « $f$ multiplié par $x$ ».
Préciser une fonction exige son domaine et sa règle ; une même formule sur deux domaines donne deux fonctions différentes.
\section{Domaine d'une expression et d'un modèle}
\begin{methode}
Pour un quotient, exclure les zéros du dénominateur.
Pour une racine carrée réelle, imposer un radicand positif ou nul.
Combiner les conditions par intersection, puis ajouter les restrictions du contexte.
\end{methode}
\begin{exemple}
$f(x)=\sqrt{x+1}/(x-2)$ exige $x\geq-1$ et $x\ne2$ :
$D=[-1,2[\cup]2,+\infty[$.
Si $A(x)=x(10-x)$ décrit l'aire d'un rectangle de périmètre $20$, le modèle exige $0<x<10$, alors que la formule existe pour tout réel.
\end{exemple}
Une fonction peut être définie par morceaux. Par exemple, un parking coûte
$P(t)=3$ pour $0<t\leq1$ heure et $P(t)=3+2(t-1)$ pour $t>1$.
Aux raccords, préciser quel intervalle inclut la borne, pour attribuer une seule sortie à chaque entrée.
\section{Courbe représentative}
\begin{definition}
La courbe de $f$ est l'ensemble des points $(x,f(x))$ pour $x\in D$.
$M(a,b)$ lui appartient si et seulement si $a\in D$ et $b=f(a)$.
\end{definition}
Une verticale rencontre cette courbe en au plus un point. Un cercle entier n'est donc pas la courbe d'une fonction de $x$, contrairement à son demi-cercle supérieur.
\begin{center}
\begin{tikzpicture}
\begin{axis}[width=10cm,height=6cm,axis lines=middle,xmin=-3,xmax=3,ymin=-.5,ymax=5,xtick={-2,-1,0,1,2},ytick={0,1,2,3,4},xlabel={$x$},ylabel={$y$},samples=100]
\addplot[accent,thick,domain=-2.2:2.2]{x^2};
\addplot[dashed,domain=-2.8:2.8]{4};
\addplot[only marks]coordinates{(-2,4)(2,4)};
\node at(axis cs:1.3,4.4){$y=4$};
\end{axis}
\end{tikzpicture}
\end{center}
Les intersections avec l'horizontale $y=4$ indiquent les antécédents de $4$.
Un graphique imprimé donne généralement des lectures approchées ; seules des coordonnées explicitement données sont exactes.
\section{Tableaux et programmes}
Un tableau échantillonne une fonction. Il ne prouve pas ce qui se passe entre deux abscisses.
\begin{exemple}
$f(x)=0$ et $g(x)=x(x-1)(x-2)$ donnent toutes deux $0$ en $0,1,2$ ; pourtant $g(1/2)=3/8\ne0$.
Un tableau n'établit donc ni une identité ni la monotonie sur un intervalle.
\end{exemple}
\begin{lstlisting}[language=Python]
def aire_rectangle(x):
    return x * (10 - x)

for k in range(1, 10):
    print(k, aire_rectangle(k))
\end{lstlisting}
Ce programme calcule les images pour les entiers de $1$ à $9$, pas pour tous les réels. Le maximum parmi ces neuf valeurs n'est pas encore une preuve du maximum global.
\section{Continu ou discret ?}
Une grandeur de temps peut se modéliser sur un intervalle réel ; un nombre d'objets se modélise sur $\N$.
Par exemple, $u(n)=12+3n$ est le coût d'envoi de $n$ objets. Les points ont des abscisses entières ; relier les points ne crée pas des demi-objets admissibles.
\begin{exemple}[Un modèle et ses limites]
$h(t)=20-5t^2$ peut modéliser une chute simplifiée à partir de $20$ mètres, avec $t$ en secondes.
On impose $t\geq0$ et $h(t)\geq0$, soit $0\leq t\leq2$.
La formule ne dit pas qu'un objet traverse le sol après $2$ secondes : le modèle cesse alors de s'appliquer.
\end{exemple}
\section{Une lecture graphique contrôlée par le calcul}
Sur $[-2;3]$, représentons $f(x)=x^2-1$ et $g(x)=x+1$.
\begin{center}
\begin{tikzpicture}
\begin{axis}[width=11cm,height=7cm,axis lines=middle,xmin=-2.4,xmax=3.4,ymin=-2,ymax=9,
xlabel={$x$},ylabel={$y$},grid=major,xtick={-2,-1,0,1,2,3},
legend style={at={(.03,.97)},anchor=north west,font=\small}]
\addplot[accent,thick,domain=-2:3,samples=101]{x^2-1};
\addlegendentry{$f(x)=x^2-1$}
\addplot[black,dashed,domain=-2:3]{x+1};
\addlegendentry{$g(x)=x+1$}
\end{axis}
\end{tikzpicture}
\end{center}
On lit $f(0)=-1$ et les antécédents de zéro $-1$ et $1$.
La courbe de f est sous l'axe des abscisses pour $-1<x<1$.
Les intersections avec g semblent avoir pour abscisses $-1$ et $2$ ; le calcul le confirme :
$f(x)-g(x)=x^2-x-2=(x+1)(x-2)$.
Le tableau de signes donne $f(x)\le g(x)$ exactement pour $x\in[-1;2]$.
Pour une lecture seulement graphique, annoncer des valeurs approchées lorsque les graduations ne permettent pas une lecture exacte.
\utilite{Modéliser explicite une dépendance et son domaine. Une bonne définition précède tout calcul de limite, de dérivée ou d'intégrale.}
\begin{retenir}
Distinguer entrée, image et antécédent ; fixer le domaine ; choisir une représentation et annoncer ce qu'elle permet de conclure.
\end{retenir}
\end{document}
